\chapter[Seguimiento del backlog]{Seguimiento del backlog en las Etapas 2 y 3}
\label{anx:seguimiento}
\small
\setlength{\tabcolsep}{4pt}

Estado de cada ítem planificado para la Etapa 2 al cierre de la etapa, más los ítems de las Etapas 3 y 4
que se adelantaron, y estado de los elementos de la Etapa 3 al 15 de octubre. La evidencia remite a las figuras
del capítulo~\ref{cap:diseno}, al capítulo~\ref{cap:pruebas} y a la documentación del repositorio, donde vive el
registro (\texttt{docs/12} y \texttt{docs/18}); un verificador comprueba que estado y horas de estos cuadros coincidan con él.

\begin{longtable}{@{}L{1.2cm}L{4.6cm}C{0.6cm}L{5.4cm}@{}}
\caption{Backlog de la Etapa 2 frente a lo ejecutado, con los ítems adelantados.}\label{tab:seguimiento}\\
\filaencab \encab{Ítem} & \encab{Tarea} & \encab{h} & \encab{Estado al cierre} \\ \midrule \endfirsthead
\filaencab \encab{Ítem} & \encab{Tarea} & \encab{h} & \encab{Estado al cierre} \\ \midrule \endhead
\multicolumn{4}{r}{\footnotesize\itshape continúa en la página siguiente} \\ \endfoot
\endlastfoot
TT-01 & Componer el entorno con contenedores: TimescaleDB y Mosquitto orquestados (ChirpStack, solo si se activa la ruta A), con un único comando de arranque & 6 & Hecho: docker-compose con TimescaleDB y Mosquitto; ChirpStack diseñado, opcional desde la Etapa 3 (ruta A) \\
TT-02 & Inicializar el proyecto TypeScript + NestJS con estructura de módulos correspondiente a M1–M8 & 4 & Hecho: backend NestJS con módulos M1, M2, M4, M5, M6 \\
TT-03 & Configurar verificación automática: tipos, formato y ejecución de pruebas & 3 & Hecho: comprobación de tipos, ESLint con reglas tipadas y pruebas unitarias de las funciones puras del pipeline; guardianes de concordancia \\
TT-04 & Documentar el procedimiento de puesta en marcha en el repositorio & 1 & Hecho: READMEs de plataforma y backend \\
TT-05 & Diagrama de arquitectura de componentes, con los ocho módulos y sus dependencias & 5 & Hecho: vista de contenedores C4 (figura~\ref{fig:contenedores}) \\
TT-06 & Diagrama de topología de red con segmentación de dispositivos IoT (RNF-09) y sentido de las conexiones (RNF-10) & 6 & Hecho: topología de red (figura~\ref{fig:red}) \\
TT-07 & Modelo de datos entidad-relación con al menos ocho entidades & 6 & Hecho: ER de 18 entidades generado desde la base (figura~\ref{fig:modelo-datos}) \\
TT-08 & Diagrama de flujo de recepción, desde la recepción hasta la persistencia, para los tres transportes & 4 & Hecho: secuencia del pipeline (figura~\ref{fig:sec-pipeline}) \\
TT-09 & Diagrama de secuencia del aprovisionamiento de un dispositivo nuevo & 3 & Hecho: secuencia de aprovisionamiento (figura~\ref{fig:sec-alta}) \\
TT-10 & Registro de decisiones de diseño con su justificación & 2 & Hecho: tablas de decisiones en el cap. 4 y bitácora \\
TT-11 & Esquema del modelo interno y contrato del tópico MQTT (nombres, unidades, identidad) & 3 & Hecho: docs/08, docs/10 §1; tópico sitio/·/dispositivo/·/magnitud \\
TT-12 & Catálogo de tipos de sensor y unidades aceptadas, con umbrales de disparo y restauración & 6 & Parcial: 7 magnitudes con rango; ampliar con los sensores reales \\
TT-13 & Especificación del contrato de la API en OpenAPI, generada desde el código & 3 & Hecho: generado desde los decoradores; servido en \texttt{/docs} y versionado (\texttt{openapi.json}) \\
TT-14 & Política de QoS, reintento y cuarentena & 2 & Hecho: QoS 1 con sesión persistente, inserción idempotente, cuarentena con motivo \\
TT-15 & Procedimiento de asignación de identidad y gestión de claves AES por sitio & 2 & Parcial: gestión de credenciales documentada (docs/11 §7); claves LoRaWAN las guarda el servidor de red (la pasarela en la ruta B, ChirpStack en la A) \\
HU-01 & Publicar una trama SenML por MQTT o HTTP desde un sensor sin infraestructura intermedia & 6 & Parcial: MQTT sí; HTTP en la Etapa 3 si queda tiempo, si no en la 4 \\
TT-16 & Validador de esquema SenML y analizador de tramas & 3 & Hecho: adaptador SenML \\
TT-19 & Esquema de base de datos y migraciones versionadas & 5 & Hecho: esquema + 16 migraciones \\
TT-24 & Comparar al menos tres alternativas de alojamiento con criterios ponderados & 4 & Hecho: cinco alternativas, seis criterios (cuadro~\ref{tab:alojamiento}) \\
TT-25 & Cuantificar el costo mensual y anual de operación en CLP y USD & 3 & Hecho: precios fechados, convertidos con dólar observado \\
TT-26 & Redactar la recomendación de despliegue justificada & 1 & Hecho: VPS de 2 vCPU/4 GB; nacional si se exigen datos en Chile \\
TT-35 & Documento de diseño de la Etapa 2, con todos los diagramas & 6 & Hecho: este informe (30 págs.), con la retroalimentación incorporada (sección~\ref{sec:retro}) \\
HU-03 & Publicar eventos en MQTT desde dispositivos que ya usan ese transporte & 4 & \textbf{Adelantado} (plan.\ E3), hecho: publicado y probado con mosquitto\_pub \\
HU-04 & Adaptador para un formato no SenML (extensibilidad de la recepción) & 7 & \textbf{Adelantado} (plan.\ E4), parcial: adaptador plano; se cierra con la trama real \\
TT-17 & Pipeline de normalización al modelo interno (seis filtros) con cuarentena & 6 & \textbf{Adelantado} (plan.\ E3), hecho: seis filtros verificados \\
HU-06 & Configurar umbrales por dispositivo para generar un evento al cruzarlos & 4 & \textbf{Adelantado} (plan.\ E4), parcial: motor listo; falta interfaz \\
TT-18 & Normalización de marcas temporales, tolerando nodos sin reloj confiable & 2 & \textbf{Adelantado} (plan.\ E3), hecho: tiempoPlausible \\
TT-20 & Autenticación de la recepción mediante credencial por sitio & 4 & \textbf{Adelantado} (plan.\ E4), hecho: credencial por dispositivo + ACL \\
\end{longtable}


\begin{longtable}{@{}L{1.2cm}L{4.6cm}C{0.6cm}L{5.4cm}@{}}
\caption{Elementos de la Etapa 3 al 15 de octubre (tras el corte del 8/10 y los adelantos del 12/10).}\label{tab:seguimiento3}\\
\filaencab \encab{Ítem} & \encab{Tarea} & \encab{h} & \encab{Estado al 15/10} \\ \midrule \endfirsthead
\filaencab \encab{Ítem} & \encab{Tarea} & \encab{h} & \encab{Estado al 15/10} \\ \midrule \endhead
\multicolumn{4}{r}{\footnotesize\itshape continúa en la página siguiente} \\ \endfoot
\endlastfoot
HU-02 & Eventos LoRaWAN en la plataforma & 6 & Parcial: adaptadores del ChirpStack del campus (ruta C) y de la pasarela Milesight (ruta B), tópico de pasarela y pasarelas por API, probados en el banco con el JSON oficial; falta la prueba en terreno y asociar RSSI y SNR (HU-15) \\
HU-03 & Publicar eventos en MQTT & 4 & Hecho: publicado y probado por MQTT \\
TT-17 & Pipeline de seis filtros con cuarentena & 6 & Hecho: seis filtros verificados \\
HU-05 & Condiciones abiertas por dispositivo & 7 & Hecho: una condición abierta por regla; pruebas de carrera \\
TT-18 & Marcas de tiempo de nodos sin reloj confiable & 2 & Hecho: marca del emisor solo si es plausible \\
HU-10 & Inventario con estado y vía de envío & 5 & Hecho: vía y formato en la API y en el panel \\
HU-11 & Registrar un dispositivo desde la plataforma & 5 & Hecho: alta, baja y rotación por API, todo o nada con el broker \\
HU-14 & Detectar que un dispositivo dejó de reportar & 7 & Hecho: supervisor con intervalo por dispositivo, medido \\
TT-39 & Canal de tiempo real (SSE) & 4 & Parcial: canal en vivo hecho; falta la sesión (Etapa 4) \\
HU-19 & Panel de salud con semáforo & 9 & Parcial: pantalla hecha; falta verla con el sensor físico \\
HU-20 & Inventario filtrable por sitio, zona, estado y vía & 7 & Parcial (adelantada): filtros por estado, zona, vía y búsqueda; falta el de sitio \\
HU-08 & Exportar los datos en CSV, XLSX y JSON & 4 & Parcial (adelantada): CSV y JSON; XLSX con la sesión \\
HU-21 & Ficha del dispositivo: condiciones, historial y serie & 7 & Parcial (adelantada): ficha con gráfico; falta el rango a medida \\
TT-28 & Emulador de nodos & 5 & Hecho: emulador de nodos para carga y pruebas \\
TT-32 & Sensor físico de extremo a extremo & 3 & Pendiente: placa recibida; sensores en tránsito \\
TT-36 & Informe de pruebas y validación preliminar & 5 & Parcial: capítulo~\ref{cap:pruebas} y documentación de pruebas \\
\end{longtable}

\section*{Retroalimentación de las Etapas 1 y 2}

\begin{table}[H]
\centering\footnotesize
\begin{tablaseg}
\begin{tabular}{@{}L{4.1cm}L{6.0cm}L{2.5cm}@{}}
\toprule
\textbf{Observación} & \textbf{Cambio realizado} & \textbf{Dónde} \\
\midrule
Diagnóstico directo con la organización destinataria pendiente; sin percepción de las partes interesadas &
Contacto obtenido por el profesor guía el 29/09: un académico de la Facultad que incorporará dispositivos IoT. Entrevista enviada por correo ese mismo día (diez preguntas: magnitudes, cantidad, uso de los datos, tiempo real, alertas, exportación), respondida el 5/10: habilita la fuente [ENT] (sección~\ref{sec:retro}) &
OE1; §\ref{sec:levantamiento}; R2 \\
Se alterna entre organizaciones de La Araucanía y la UCT como destinatario &
El problema se centra solo en la UCT, con el campus como caso de validación; ninguna otra organización ha manifestado interés y la generalización queda fuera del alcance &
§\ref{sec:problema}; resumen \\
Causas respaldadas solo con la oferta comercial &
Párrafo de evidencia con literatura sobre interoperabilidad y dependencia de proveedor (C1 a C3); la carga técnica de las alternativas (C4) pasa al estado del arte y los costos al anexo~\ref{anx:evaluacion}; la evidencia directa se asigna al levantamiento &
§\ref{sec:problema}; §\ref{sec:estado-arte} \\
OE1 fija el levantamiento al 4 de septiembre pero se declara pendiente &
OE1 con plazo al 30 de octubre; la evidencia recogida antes del 4 de septiembre (indicaciones del profesor guía, mercado, base teórica) queda explicitada en §\ref{sec:levantamiento} y el catálogo se valida con las entrevistas &
§\ref{sec:objetivos} \\
Alcance funcional comprometido demasiado amplio &
Criterio de demostración en dos niveles (§\ref{sec:backlog}): el núcleo mínimo demostrable primero, el resto contra holgura; OE4 reescrito en ese orden &
OE4; §\ref{sec:backlog} \\
OE5 sin umbrales de validación &
OE5 con umbrales explícitos —casos de prueba, carga, latencia, detección de caída, recuperación, disponibilidad, memoria— tomados de los requerimientos no funcionales &
OE5 \\
\bottomrule
\end{tabular}
\end{tablaseg}
\caption{Observaciones del profesor guía sobre la Etapa 1 y cambios derivados en la Etapa 2.}
\label{tab:retro}
\end{table}

\begin{table}[H]
\centering\footnotesize
\begin{tablaseg}
\begin{tabular}{@{}L{4.1cm}L{6.0cm}L{2.5cm}@{}}
\toprule
\textbf{Observación (29/09)} & \textbf{Cambio realizado} & \textbf{Dónde} \\
\midrule
«Plataforma web» se lee como una página; es un sistema integrado &
Título: \emph{Sistema integrado para el registro, gestión y monitoreo de dispositivos IoT} &
Portada \\
Objetivo general largo y con detalle técnico (TCP/IP) &
Reescrito sin detalle técnico y con el para quién (investigación de la Facultad); 82 a 47 palabras &
§\ref{sec:objetivos} \\
Objetivos específicos extensos y dispersos &
Los cinco conservan número y trazabilidad, 23\,\% más cortos; OE4 suma el registro de dispositivos y OE5 mantiene sus umbrales &
§\ref{sec:objetivos} \\
El problema debe entenderse desde el primer párrafo &
La sección abre con el problema y un caso verificado; sin detalle de hardware ni protocolos; alcance desde el registro hasta el dashboard &
§\ref{sec:problema} \\
Estado del arte poco claro; no se identifican los competidores &
Tres grupos con afirmaciones citadas de fuentes oficiales, y la ubicación de este trabajo frente a ellos &
§\ref{sec:estado-arte} \\
Resumen extenso y poco claro &
Reescrito en dos párrafos: problema, propuesta, método, resultado y siguiente etapa; 344 a 263 palabras &
Resumen \\
\bottomrule
\end{tabular}
\end{tablaseg}
\caption{Observaciones del profesor guía en la reunión del 29 de septiembre y cambios derivados.}
\label{tab:retro2}
\end{table}

\section*{Calendario oficial}

\begin{table}[H]
\centering\small
\begin{tabular}{@{}L{6.6cm}C{2.8cm}C{3.0cm}@{}}
\toprule
\textbf{Etapa} & \textbf{Sprint} & \textbf{Entrega} \\
\midrule
1 --- Problema, objetivos, requerimientos & 17 ago -- 4 sep & \textbf{4 de septiembre} \\
2 --- Diseño e implementación inicial & 7 sep -- 2 oct & 2 de octubre \\
3 --- Mejoras, pruebas, validación preliminar & 5 -- 30 oct & 30 de octubre \\
4 --- Validación final e informe & 2 nov -- 4 dic & 4 de diciembre \\
\bottomrule
\end{tabular}
\caption{Calendario oficial (pauta del 17 de agosto de 2026). Presupuesto a $\sim$15 h semanales:
$\sim$240 horas.}
\label{tab:calendario}
\end{table}


\section*{Registro de riesgos}

\begin{table}[H]
\centering\small
\begin{tabular}{@{}lL{3.8cm}L{2.2cm}L{5.0cm}@{}}
\toprule
 & \textbf{Riesgo} & \textbf{Estado} & \textbf{Mitigación} \\
\midrule
R1 & Hardware de sensores no disponible & \textbf{En tratamiento} & Emulador de nodo por software para la carga; para el sensor real (TT-32, núcleo), plan B activado el 7/10: nodo Wi-Fi propio, con la placa recibida el 12/10 y los sensores en tránsito \\
R2 & Sin acceso a partes interesadas & \textbf{En tratamiento} & Entrevista respondida el 5/10 (sección~\ref{sec:retro}); falta la segunda (OE1) \\
R3 & Ambigüedad de fechas & \textbf{Cerrado} & Calendario oficial publicado el 24/08 \\
R4 & Alcance excesivo & \textbf{Confirmado y controlado} & Núcleo comprometido vs.\ alcance condicional; revisión de holgura por etapa \\
R5 & Protocolo inadecuado al patrón de uso & \textbf{Cerrado} & Se materializó y resolvió: Contact ID descartado el 02/09; el transporte es MQTT y el formato SenML (sección~\ref{sec:transporte}) \\
R6 & Recepción LoRaWAN en terreno sin una pasarela que alcance el servidor & \textbf{En tratamiento} & Banco de pruebas con el JSON oficial; prueba en la misma red local que la pasarela del campus; servidor alcanzable desde el campus, por gestionar \\
\bottomrule
\end{tabular}
\caption{Registro de riesgos.}
\label{tab:riesgos}
\end{table}
